% Источники: 07_Программные_требования/DOC-06_Программные_требования.md
% Генератор: docs/build_pdf.py
% Шаблон печатного представления. Содержимое подставляет docs/build_pdf.py.
\documentclass[12pt,a4paper]{article}
\usepackage{iftex}
\ifPDFTeX
  \usepackage[utf8]{inputenc}
  \usepackage[T2A]{fontenc}
  \usepackage{newtxtext,newtxmath}
\else
  \usepackage{fontspec}
  \IfFontExistsTF{Times New Roman}
    {\setmainfont{Times New Roman}}
    {\setmainfont{Libertinus Serif}}
  \IfFontExistsTF{Arial}{\setsansfont{Arial}}{\setsansfont{Latin Modern Sans}}
  \IfFontExistsTF{Menlo}{\setmonofont{Menlo}[Scale=MatchLowercase]}
    {\setmonofont{Latin Modern Mono}[Scale=MatchLowercase]}
\fi
\usepackage[russian]{babel}
\usepackage{geometry}
\geometry{left=25mm,right=15mm,top=20mm,bottom=20mm}
\usepackage{setspace}
\onehalfspacing
\usepackage{indentfirst}
\setlength{\parindent}{1.25cm}
\setlength{\parskip}{2pt}
\setlength{\emergencystretch}{3em}
\usepackage{graphicx,booktabs,longtable,array,tabularx,calc,etoolbox}
\usepackage{enumitem}
\setlist[itemize]{leftmargin=1.25cm,itemsep=2pt,topsep=4pt}
\setlist[enumerate]{leftmargin=1.25cm,itemsep=2pt,topsep=4pt}
\newcommand{\tightlist}{\setlength{\itemsep}{0pt}\setlength{\parskip}{0pt}}
\AtBeginEnvironment{longtable}{\small\singlespacing}
\setlength{\LTpre}{8pt}
\setlength{\LTpost}{8pt}
\setlength{\tabcolsep}{5pt}
\setlength{\LTleft}{0pt}
\setlength{\LTright}{0pt}
\usepackage{xurl}
\usepackage[unicode,colorlinks=true,linkcolor=black,urlcolor=blue]{hyperref}
\usepackage{bookmark}
\hypersetup{pdftitle={Сводные программные требования},pdfauthor={Анохин Г.А.},pdfsubject={Криптотренажёр mephi-12}}
\urlstyle{same}
\pagestyle{plain}
\setcounter{secnumdepth}{3}
\setcounter{tocdepth}{2}

\begin{document}
\hypersetup{pageanchor=false}
\begin{titlepage}
  \thispagestyle{empty}
  \centering
  {\large НИЯУ МИФИ\par}
  \vspace*{3.5cm}
  {\LARGE\bfseries Сводные программные требования\par}
  \vspace{0.5cm}
  {\large DOC-06\par}
  \vspace{0.6cm}
  {\large Платформа для изучения и практики курса по криптографии\par}
  {\large Криптотренажёр mephi-12\par}
  \vspace{0.7cm}
  {\large Версия 0.6\par}
  {\large Редакция от 07.11.2026\par}
  \vspace{1cm}
  \begin{tabularx}{\textwidth}{@{}p{3.4cm}Xp{4cm}@{}}
    \toprule
    Роль & ФИО & Подпись / Дата \\
    \midrule
    Разработал & Анохин Г.А. & \\
    \addlinespace[6pt]
    Согласовано & & \\
    \addlinespace[6pt]
    Утверждено & & \\
    \bottomrule
  \end{tabularx}
  \vfill
  {\large Москва, 2026\par}
\end{titlepage}
\hypersetup{pageanchor=true}

\section*{Лист регистрации изменений}
\addcontentsline{toc}{section}{Лист регистрации изменений}
\begingroup\small\singlespacing
\noindent
\begin{tabularx}{\textwidth}{@{}p{1.2cm}p{2cm}p{2.5cm}Xp{1.8cm}@{}}
\toprule
Версия & Дата & Автор & Содержание изменения & Утвердил \\
\midrule
0.6 & 07.11.2026 & Анохин Г.А. & Уточнены правила и контракты; расширено описание сценариев; подготовлена печатная редакция. & \\
\bottomrule
\end{tabularx}
\endgroup
\clearpage
\begingroup\small\singlespacing
\tableofcontents
\endgroup
\clearpage

\hypertarget{doc-06-сводные-программные-требования}{}

{\def\LTcaptype{none} % do not increment counter
\begin{longtable}[]{@{}
  >{\raggedright\arraybackslash}p{(\linewidth - 2\tabcolsep) * \real{0.3200}}
  >{\raggedright\arraybackslash}p{(\linewidth - 2\tabcolsep) * \real{0.6800}}@{}}
\toprule\noalign{}
\begin{minipage}[b]{\linewidth}\raggedright
Поле
\end{minipage} & \begin{minipage}[b]{\linewidth}\raggedright
Значение
\end{minipage} \\
\midrule\noalign{}
\endhead
\bottomrule\noalign{}
\endlastfoot
Версия & 0.6 \\
Дата & 07.11.2026 \\
Редактор & Анохин Г.А. \\
Основание & Рабочая редакция DOC-05 \\
\end{longtable}
}

\section{Назначение и границы}\label{1-ux43dux430ux437ux43dux430ux447ux435ux43dux438ux435-ux438-ux433ux440ux430ux43dux438ux446ux44b}

\hypertarget{1-назначение-и-границы}{}

Документ распределяет общие системные обязательства между логическими программными компонентами и определяет условия их совместной работы. Подсистемные требования уточняют эти обязательства; в текущую редакцию включён \href{DOC-06.2_\%D0\%A1\%D0\%B5\%D1\%81\%D1\%81\%D0\%B8\%D0\%B8_\%D0\%B7\%D0\%B0\%D0\%B4\%D0\%B0\%D1\%87_\%D0\%B8_\%D0\%BF\%D0\%BB\%D0\%B0\%D0\%BD\%D0\%B8\%D1\%80\%D0\%BE\%D0\%B2\%D1\%89\%D0\%B8\%D0\%BA.pdf\#nameddest=doc-062-\%D1\%81\%D0\%B5\%D1\%81\%D1\%81\%D0\%B8\%D0\%B8-\%D0\%B7\%D0\%B0\%D0\%B4\%D0\%B0\%D1\%87-\%D0\%B8-\%D0\%BF\%D0\%BB\%D0\%B0\%D0\%BD\%D0\%B8\%D1\%80\%D0\%BE\%D0\%B2\%D1\%89\%D0\%B8\%D0\%BA}{DOC-06.2 по сессиям}.

Компонент здесь означает область ответственности ПО. Такое разделение не требует отдельных сервисов или применения единственного архитектурного стиля. Математические алгоритмы, модель хранения и полный справочник HTTP-операций не дублируются в общей части.

\section{Ответственность и связи}\label{2-ux43eux442ux432ux435ux442ux441ux442ux432ux435ux43dux43dux43eux441ux442ux44c-ux438-ux441ux432ux44fux437ux438}

\hypertarget{2-ответственность-и-связи}{}

{\def\LTcaptype{none} % do not increment counter
\begin{longtable}[]{@{}
  >{\raggedright\arraybackslash}p{(\linewidth - 4\tabcolsep) * \real{0.2468}}
  >{\raggedright\arraybackslash}p{(\linewidth - 4\tabcolsep) * \real{0.4774}}
  >{\raggedright\arraybackslash}p{(\linewidth - 4\tabcolsep) * \real{0.2759}}@{}}
\toprule\noalign{}
\begin{minipage}[b]{\linewidth}\raggedright
Компонент
\end{minipage} & \begin{minipage}[b]{\linewidth}\raggedright
Вход и результат
\end{minipage} & \begin{minipage}[b]{\linewidth}\raggedright
Ответственный за предметную часть
\end{minipage} \\
\midrule\noalign{}
\endhead
\bottomrule\noalign{}
\endlastfoot
Проверка доступа & Учётные данные запроса; подтверждённый пользователь и полномочия либо отказ & Кокорев \\
Конфигурации & Пользователь и его группа; доступный состав набора либо отказ & Самбьюи \\
Управление сессиями & Запрос создания или ответа; сессия и предметный результат либо категория отказа & Анохин \\
Генераторы и проверяющие & Тип и настройки; условие. Сохранённое условие и ответ; предметный вердикт либо ошибка проверки & Владельцы криптографических доменов \\
Контроль срока & Сохранённые времена и сведения о принятых ответах; завершение просроченных задач & Анохин \\
Хранение & Данные и изменения состояния; согласованное сохранение и чтение & Анохин для правил сессий, Самбьюи для общей модели и инфраструктуры \\
\end{longtable}
}

При создании набора управление сессиями получает пользователя и доступную конфигурацию, вызывает генераторы и сохраняет состав. При отправке ответа оно проверяет допуск, обращается к профильному проверяющему и сохраняет результат. Контроль срока использует те же данные времени и завершения; отдельного правила подсчёта итога у него нет.

\subsection{Последовательность создания набора}\label{ux43fux43eux441ux43bux435ux434ux43eux432ux430ux442ux435ux43bux44cux43dux43eux441ux442ux44c-ux441ux43eux437ux434ux430ux43dux438ux44f-ux43dux430ux431ux43eux440ux430}

\hypertarget{последовательность-создания-набора}{}

Компонент приёма операции передаёт управлению сессиями подтверждённого пользователя и имя конфигурации. Затем определяется группа пользователя и запрашивается доступный состав набора. Эти действия не заменяются доверием к переданному клиентом идентификатору группы.

По составу вызываются соответствующие генераторы. Полученные условия, типы и принадлежность задач образуют данные создаваемой сессии. Подтверждение выдаётся после сохранения полного набора. При отказе генерации или сохранения частичный набор не должен оставаться доступной активной сессией.

Ограничение одной активной сессии охватывает всю операцию создания. Отдельная предварительная проверка без согласования последующего сохранения недостаточна: два обращения могли бы одновременно увидеть отсутствие сессии. Конкретный механизм защиты выбирается при проектировании хранения, но программное обязательство уже определено.

\subsection{Последовательность обработки ответа}\label{ux43fux43eux441ux43bux435ux434ux43eux432ux430ux442ux435ux43bux44cux43dux43eux441ux442ux44c-ux43eux431ux440ux430ux431ux43eux442ux43aux438-ux43eux442ux432ux435ux442ux430}

\hypertarget{последовательность-обработки-ответа}{}

Приём запроса фиксирует серверное время его поступления. Управление сессиями устанавливает задачу, проверяет принадлежность, состояние и срок, затем проверяет формат по схеме сохранённого типа. Предметный обработчик получает только допустимый ответ с условием выбранной задачи.

После получения вердикта управление сессиями согласованно сохраняет результат задачи и вычисляет итог набора. Контроль срока использует те же правила завершения и сведения о своевременном запросе. Если ответ получен до срока, момент окончания вычисления не отменяет его своевременность.

Общие требования к обмену определяют смысл данных на этих переходах. Поля и HTTP-коды раскрываются в интерфейсных спецификациях; математические критерии правильности --- в предметных подсистемах. Такое распределение позволяет уточнять внутреннее устройство части, сохраняя её обязательства перед остальными участниками.

\section{Общие программные требования}\label{3-ux43eux431ux449ux438ux435-ux43fux440ux43eux433ux440ux430ux43cux43cux43dux44bux435-ux442ux440ux435ux431ux43eux432ux430ux43dux438ux44f}

\hypertarget{3-общие-программные-требования}{}

Все требования ниже обязательны. Ссылки SR ведут к \href{../04_\%D0\%A1\%D0\%B8\%D1\%81\%D1\%82\%D0\%B5\%D0\%BC\%D0\%BD\%D1\%8B\%D0\%B5_\%D1\%82\%D1\%80\%D0\%B5\%D0\%B1\%D0\%BE\%D0\%B2\%D0\%B0\%D0\%BD\%D0\%B8\%D1\%8F/DOC-05_\%D0\%A1\%D0\%B8\%D1\%81\%D1\%82\%D0\%B5\%D0\%BC\%D0\%BD\%D1\%8B\%D0\%B5_\%D1\%82\%D1\%80\%D0\%B5\%D0\%B1\%D0\%BE\%D0\%B2\%D0\%B0\%D0\%BD\%D0\%B8\%D1\%8F.pdf\#nameddest=doc-05-\%D1\%82\%D1\%80\%D0\%B5\%D0\%B1\%D0\%BE\%D0\%B2\%D0\%B0\%D0\%BD\%D0\%B8\%D1\%8F-\%D0\%BA-\%D0\%B8\%D0\%BD\%D1\%82\%D0\%B5\%D0\%B3\%D1\%80\%D0\%B0\%D1\%86\%D0\%B8\%D0\%B8}{системному разделу интеграции}; проверки являются описательными.

\subsection{SWR-SYS-001. Передача личности пользователя}\label{swr-sys-001-ux43fux435ux440ux435ux434ux430ux447ux430-ux43bux438ux447ux43dux43eux441ux442ux438-ux43fux43eux43bux44cux437ux43eux432ux430ux442ux435ux43bux44f}

\hypertarget{swr-sys-001-передача-личности-пользователя}{}

Входной компонент проверяет доступ и передаёт остальным частям подтверждённого пользователя. Дальнейшие операции используют этот результат как контекст выполнения. Имя конфигурации, UUID задачи и другие параметры остаются данными операции; они не назначают ей другого владельца. Системное основание такого разделения дано в \href{../04_\%D0\%A1\%D0\%B8\%D1\%81\%D1\%82\%D0\%B5\%D0\%BC\%D0\%BD\%D1\%8B\%D0\%B5_\%D1\%82\%D1\%80\%D0\%B5\%D0\%B1\%D0\%BE\%D0\%B2\%D0\%B0\%D0\%BD\%D0\%B8\%D1\%8F/DOC-05_\%D0\%A1\%D0\%B8\%D1\%81\%D1\%82\%D0\%B5\%D0\%BC\%D0\%BD\%D1\%8B\%D0\%B5_\%D1\%82\%D1\%80\%D0\%B5\%D0\%B1\%D0\%BE\%D0\%B2\%D0\%B0\%D0\%BD\%D0\%B8\%D1\%8F.pdf\#nameddest=sr-sys-001-\%D0\%B5\%D0\%B4\%D0\%B8\%D0\%BD\%D0\%B0\%D1\%8F-\%D0\%BB\%D0\%B8\%D1\%87\%D0\%BD\%D0\%BE\%D1\%81\%D1\%82\%D1\%8C-\%D0\%BF\%D0\%BE\%D0\%BB\%D1\%8C\%D0\%B7\%D0\%BE\%D0\%B2\%D0\%B0\%D1\%82\%D0\%B5\%D0\%BB\%D1\%8F}{единой личности пользователя}.

\textbf{Требование.} Компоненты пользовательских операций должны получать личность пользователя из результата проверки доступа. Параметры тела и строки запроса не должны заменять подтверждённого владельца операции.

При создании сессии контекст используется для определения группы, при ответе --- для проверки владельца задачи, при чтении --- для отбора собственной сессии. Повторять логику выдачи полномочий в каждом предметном проверяющем не требуется. Порядок пользовательских операций определён в \href{DOC-06.2_\%D0\%A1\%D0\%B5\%D1\%81\%D1\%81\%D0\%B8\%D0\%B8_\%D0\%B7\%D0\%B0\%D0\%B4\%D0\%B0\%D1\%87_\%D0\%B8_\%D0\%BF\%D0\%BB\%D0\%B0\%D0\%BD\%D0\%B8\%D1\%80\%D0\%BE\%D0\%B2\%D1\%89\%D0\%B8\%D0\%BA.pdf\#nameddest=3-\%D0\%BE\%D0\%B1\%D1\%80\%D0\%B0\%D0\%B1\%D0\%BE\%D1\%82\%D0\%BA\%D0\%B0-\%D0\%BE\%D1\%82\%D0\%B2\%D0\%B5\%D1\%82\%D0\%B0}{обработке ответа}, а доступ по HTTP --- в \href{DOC-06.2_\%D0\%A1\%D0\%B5\%D1\%81\%D1\%81\%D0\%B8\%D0\%B8_\%D0\%B7\%D0\%B0\%D0\%B4\%D0\%B0\%D1\%87_\%D0\%B8_\%D0\%BF\%D0\%BB\%D0\%B0\%D0\%BD\%D0\%B8\%D1\%80\%D0\%BE\%D0\%B2\%D1\%89\%D0\%B8\%D0\%BA.pdf\#nameddest=\%D0\%BE\%D0\%BF\%D0\%B5\%D1\%80\%D0\%B0\%D1\%86\%D0\%B8\%D0\%B8}{контракте операций}.

Источник: SR-SYS-001. Проверка: подстановка чужого идентификатора не изменяет пользователя, от имени которого выполняется операция.

\subsection{SWR-SYS-002. Общие идентификаторы}\label{swr-sys-002-ux43eux431ux449ux438ux435-ux438ux434ux435ux43dux442ux438ux444ux438ux43aux430ux442ux43eux440ux44b}

\hypertarget{swr-sys-002-общие-идентификаторы}{}

Создание, хранение, проверка и выдача результата обращаются к одному набору объектов. Программные компоненты должны передавать уже назначенные ID, а не создавать похожие идентификаторы при каждом преобразовании данных. Это особенно важно при формировании ответов и восстановлении после перезапуска, когда объект проходит через несколько представлений.

\textbf{Требование.} Компоненты должны сохранять и передавать идентификаторы пользователя, сессии и задачи без повторного назначения при обмене между ними.

ID задачи в выдаче используется как \texttt{statementId}, а результат содержит этот же ID как \texttt{problemId}. По связи с сессией определяется владелец. Формат передачи не должен менять принадлежность или связывать результат с соседней задачей одного типа. Поля описаны в \href{DOC-06.2_\%D0\%A1\%D0\%B5\%D1\%81\%D1\%81\%D0\%B8\%D0\%B8_\%D0\%B7\%D0\%B0\%D0\%B4\%D0\%B0\%D1\%87_\%D0\%B8_\%D0\%BF\%D0\%BB\%D0\%B0\%D0\%BD\%D0\%B8\%D1\%80\%D0\%BE\%D0\%B2\%D1\%89\%D0\%B8\%D0\%BA.pdf\#nameddest=\%D1\%80\%D0\%B5\%D0\%B7\%D1\%83\%D0\%BB\%D1\%8C\%D1\%82\%D0\%B0\%D1\%82-\%D0\%B8-\%D0\%BE\%D1\%82\%D0\%BA\%D0\%B0\%D0\%B7\%D1\%8B}{контракте результата}, а системное обязательство --- в \href{../04_\%D0\%A1\%D0\%B8\%D1\%81\%D1\%82\%D0\%B5\%D0\%BC\%D0\%BD\%D1\%8B\%D0\%B5_\%D1\%82\%D1\%80\%D0\%B5\%D0\%B1\%D0\%BE\%D0\%B2\%D0\%B0\%D0\%BD\%D0\%B8\%D1\%8F/DOC-05_\%D0\%A1\%D0\%B8\%D1\%81\%D1\%82\%D0\%B5\%D0\%BC\%D0\%BD\%D1\%8B\%D0\%B5_\%D1\%82\%D1\%80\%D0\%B5\%D0\%B1\%D0\%BE\%D0\%B2\%D0\%B0\%D0\%BD\%D0\%B8\%D1\%8F.pdf\#nameddest=sr-sys-002-\%D1\%81\%D0\%B2\%D1\%8F\%D0\%B7\%D0\%BD\%D0\%BE\%D1\%81\%D1\%82\%D1\%8C-\%D0\%B4\%D0\%B0\%D0\%BD\%D0\%BD\%D1\%8B\%D1\%85-\%D0\%B7\%D0\%B0\%D0\%B4\%D0\%B0\%D1\%87\%D0\%B8}{связности данных задачи}.

Источник: SR-SYS-002. Проверка: задача в выдаче, проверке и хранилище имеет одну принадлежность и один идентификатор.

\subsection{SWR-SYS-003. Состав созданной сессии}\label{swr-sys-003-ux441ux43eux441ux442ux430ux432-ux441ux43eux437ux434ux430ux43dux43dux43eux439-ux441ux435ux441ux441ux438ux438}

\hypertarget{swr-sys-003-состав-созданной-сессии}{}

Компонент конфигураций предоставляет доступный состав нового набора. После создания управление сессиями работает с сохранёнными задачами. Повторное получение текущей конфигурации во время проверки не должно менять ранее выданное условие. Поэтому данные выбора и данные уже созданного задания имеют разные роли в программе.

\textbf{Требование.} Управление сессиями должно сохранять состав и условия выданных задач по результату выбора доступной конфигурации; последующая проверка должна использовать эти сохранённые данные.

Сохранённые тип и условие передаются проверяющему без повторной генерации. Так ответ сопоставляется с тем заданием, которое получил пользователь. Порядок сохранения полного набора дан в \href{DOC-06.2_\%D0\%A1\%D0\%B5\%D1\%81\%D1\%81\%D0\%B8\%D0\%B8_\%D0\%B7\%D0\%B0\%D0\%B4\%D0\%B0\%D1\%87_\%D0\%B8_\%D0\%BF\%D0\%BB\%D0\%B0\%D0\%BD\%D0\%B8\%D1\%80\%D0\%BE\%D0\%B2\%D1\%89\%D0\%B8\%D0\%BA.pdf\#nameddest=2-\%D1\%81\%D0\%BE\%D0\%B7\%D0\%B4\%D0\%B0\%D0\%BD\%D0\%B8\%D0\%B5-\%D1\%81\%D0\%B5\%D1\%81\%D1\%81\%D0\%B8\%D0\%B8}{создании сессии}, а последствия ошибки --- в \href{DOC-06.2_\%D0\%A1\%D0\%B5\%D1\%81\%D1\%81\%D0\%B8\%D0\%B8_\%D0\%B7\%D0\%B0\%D0\%B4\%D0\%B0\%D1\%87_\%D0\%B8_\%D0\%BF\%D0\%BB\%D0\%B0\%D0\%BD\%D0\%B8\%D1\%80\%D0\%BE\%D0\%B2\%D1\%89\%D0\%B8\%D0\%BA.pdf\#nameddest=swr-ses-005-\%D1\%81\%D0\%BE\%D1\%85\%D1\%80\%D0\%B0\%D0\%BD\%D0\%B5\%D0\%BD\%D0\%B8\%D0\%B5-\%D1\%81\%D0\%BE\%D0\%B7\%D0\%B4\%D0\%B0\%D0\%BD\%D0\%BD\%D0\%BE\%D0\%B3\%D0\%BE-\%D0\%BD\%D0\%B0\%D0\%B1\%D0\%BE\%D1\%80\%D0\%B0}{согласованном сохранении}. Системный смысл разделения раскрыт в \href{../04_\%D0\%A1\%D0\%B8\%D1\%81\%D1\%82\%D0\%B5\%D0\%BC\%D0\%BD\%D1\%8B\%D0\%B5_\%D1\%82\%D1\%80\%D0\%B5\%D0\%B1\%D0\%BE\%D0\%B2\%D0\%B0\%D0\%BD\%D0\%B8\%D1\%8F/DOC-05_\%D0\%A1\%D0\%B8\%D1\%81\%D1\%82\%D0\%B5\%D0\%BC\%D0\%BD\%D1\%8B\%D0\%B5_\%D1\%82\%D1\%80\%D0\%B5\%D0\%B1\%D0\%BE\%D0\%B2\%D0\%B0\%D0\%BD\%D0\%B8\%D1\%8F.pdf\#nameddest=sr-sys-003-\%D1\%81\%D0\%BE\%D0\%B3\%D0\%BB\%D0\%B0\%D1\%81\%D0\%BE\%D0\%B2\%D0\%B0\%D0\%BD\%D0\%BD\%D1\%8B\%D0\%B9-\%D1\%81\%D0\%BE\%D1\%81\%D1\%82\%D0\%B0\%D0\%B2-\%D0\%BD\%D0\%B0\%D0\%B1\%D0\%BE\%D1\%80\%D0\%B0}{согласованном составе набора}.

Источники: SR-SYS-003, SR-SYS-004. Проверка: изменение конфигурации после создания не меняет условия уже выданной задачи.

\subsection{SWR-SYS-004. Контракт предметного обработчика}\label{swr-sys-004-ux43aux43eux43dux442ux440ux430ux43aux442-ux43fux440ux435ux434ux43cux435ux442ux43dux43eux433ux43e-ux43eux431ux440ux430ux431ux43eux442ux447ux438ux43aux430}

\hypertarget{swr-sys-004-контракт-предметного-обработчика}{}

У генератора и проверяющего разные входы и результаты. Генерация создаёт условие по составу набора. Проверка получает уже сохранённое условие и ответ пользователя. Её результат должен отличать правильность решения от невозможности выполнить проверку: это определяет, возникает ли предметный вердикт и расходуется ли попытка.

\textbf{Требование.} Предметный обработчик должен различать генерацию условия и проверку ответа, а результат проверки --- завершённый вердикт, ошибку формата или технический сбой без вердикта.

Проверяющий не должен самостоятельно назначать владельца, менять общий срок или закрывать всю сессию. Эти решения остаются у управления сессиями. Предметная часть отвечает за свой алгоритм и диагностику неверного расчёта. Подробное разделение видов ответа дано в \href{DOC-06.2_\%D0\%A1\%D0\%B5\%D1\%81\%D1\%81\%D0\%B8\%D0\%B8_\%D0\%B7\%D0\%B0\%D0\%B4\%D0\%B0\%D1\%87_\%D0\%B8_\%D0\%BF\%D0\%BB\%D0\%B0\%D0\%BD\%D0\%B8\%D1\%80\%D0\%BE\%D0\%B2\%D1\%89\%D0\%B8\%D0\%BA.pdf\#nameddest=\%D0\%BF\%D1\%80\%D0\%B5\%D0\%B4\%D0\%BC\%D0\%B5\%D1\%82\%D0\%BD\%D1\%8B\%D0\%B9-\%D0\%BE\%D1\%82\%D0\%B2\%D0\%B5\%D1\%82}{предметной схеме} и \href{DOC-06.2_\%D0\%A1\%D0\%B5\%D1\%81\%D1\%81\%D0\%B8\%D0\%B8_\%D0\%B7\%D0\%B0\%D0\%B4\%D0\%B0\%D1\%87_\%D0\%B8_\%D0\%BF\%D0\%BB\%D0\%B0\%D0\%BD\%D0\%B8\%D1\%80\%D0\%BE\%D0\%B2\%D1\%89\%D0\%B8\%D0\%BA.pdf\#nameddest=\%D1\%80\%D0\%B5\%D0\%B7\%D1\%83\%D0\%BB\%D1\%8C\%D1\%82\%D0\%B0\%D1\%82-\%D0\%B8-\%D0\%BE\%D1\%82\%D0\%BA\%D0\%B0\%D0\%B7\%D1\%8B}{результате операции}, а выбор обработчика --- в \href{DOC-06.2_\%D0\%A1\%D0\%B5\%D1\%81\%D1\%81\%D0\%B8\%D0\%B8_\%D0\%B7\%D0\%B0\%D0\%B4\%D0\%B0\%D1\%87_\%D0\%B8_\%D0\%BF\%D0\%BB\%D0\%B0\%D0\%BD\%D0\%B8\%D1\%80\%D0\%BE\%D0\%B2\%D1\%89\%D0\%B8\%D0\%BA.pdf\#nameddest=swr-ses-010-\%D0\%BF\%D0\%B5\%D1\%80\%D0\%B5\%D0\%B4\%D0\%B0\%D1\%87\%D0\%B0-\%D0\%BF\%D1\%80\%D0\%BE\%D1\%84\%D0\%B8\%D0\%BB\%D1\%8C\%D0\%BD\%D0\%BE\%D0\%B9-\%D0\%BF\%D1\%80\%D0\%BE\%D0\%B2\%D0\%B5\%D1\%80\%D0\%BA\%D0\%B5}{подсистемном требовании}.

Источники: SR-SYS-004, SR-SYS-009. Проверка: обработчик вызывается по сохранённому типу; сбой не возвращается как неверное решение. Форматы своих данных подтверждает владелец домена.

\subsection{SWR-SYS-005. Согласованное сохранение}\label{swr-sys-005-ux441ux43eux433ux43bux430ux441ux43eux432ux430ux43dux43dux43eux435-ux441ux43eux445ux440ux430ux43dux435ux43dux438ux435}

\hypertarget{swr-sys-005-согласованное-сохранение}{}

Сохранение связывает результат одной задачи с агрегированным состоянием её набора. Если эти изменения подтвердить независимо, чтение может получить несовместимые сведения: задача завершена, а итог сессии вычислен без неё. Программная операция должна иметь общую границу фиксации, соответствующую \href{../04_\%D0\%A1\%D0\%B8\%D1\%81\%D1\%82\%D0\%B5\%D0\%BC\%D0\%BD\%D1\%8B\%D0\%B5_\%D1\%82\%D1\%80\%D0\%B5\%D0\%B1\%D0\%BE\%D0\%B2\%D0\%B0\%D0\%BD\%D0\%B8\%D1\%8F/DOC-05_\%D0\%A1\%D0\%B8\%D1\%81\%D1\%82\%D0\%B5\%D0\%BC\%D0\%BD\%D1\%8B\%D0\%B5_\%D1\%82\%D1\%80\%D0\%B5\%D0\%B1\%D0\%BE\%D0\%B2\%D0\%B0\%D0\%BD\%D0\%B8\%D1\%8F.pdf\#nameddest=sr-sys-005-\%D1\%81\%D0\%BE\%D0\%B3\%D0\%BB\%D0\%B0\%D1\%81\%D0\%BE\%D0\%B2\%D0\%B0\%D0\%BD\%D0\%BD\%D0\%BE\%D1\%81\%D1\%82\%D1\%8C-\%D1\%80\%D0\%B5\%D0\%B7\%D1\%83\%D0\%BB\%D1\%8C\%D1\%82\%D0\%B0\%D1\%82\%D0\%B0-\%D0\%B8-\%D1\%81\%D0\%BE\%D1\%81\%D1\%82\%D0\%BE\%D1\%8F\%D0\%BD\%D0\%B8\%D1\%8F}{системному результату сохранения}.

\textbf{Требование.} Изменение состояния задачи и вычисление соответствующего состояния сессии должны сохраняться как одна согласованная операция до выдачи подтверждённого результата.

Математическое вычисление выполняется до короткой транзакции записи. В транзакции проверяется допустимость фиксации, сохраняется вердикт и пересчитывается набор. После её завершения ответ может подтверждать результат клиенту. Порядок блокировки и конкуренция нескольких ответов раскрыты в \href{DOC-06.2_\%D0\%A1\%D0\%B5\%D1\%81\%D1\%81\%D0\%B8\%D0\%B8_\%D0\%B7\%D0\%B0\%D0\%B4\%D0\%B0\%D1\%87_\%D0\%B8_\%D0\%BF\%D0\%BB\%D0\%B0\%D0\%BD\%D0\%B8\%D1\%80\%D0\%BE\%D0\%B2\%D1\%89\%D0\%B8\%D0\%BA.pdf\#nameddest=\%D1\%81\%D0\%BE\%D0\%B3\%D0\%BB\%D0\%B0\%D1\%81\%D0\%BE\%D0\%B2\%D0\%B0\%D0\%BD\%D0\%B8\%D0\%B5-\%D0\%BF\%D0\%B0\%D1\%80\%D0\%B0\%D0\%BB\%D0\%BB\%D0\%B5\%D0\%BB\%D1\%8C\%D0\%BD\%D1\%8B\%D1\%85-\%D1\%81\%D0\%BE\%D0\%B1\%D1\%8B\%D1\%82\%D0\%B8\%D0\%B9}{согласовании параллельных событий}; восстановление использует \href{DOC-06.2_\%D0\%A1\%D0\%B5\%D1\%81\%D1\%81\%D0\%B8\%D0\%B8_\%D0\%B7\%D0\%B0\%D0\%B4\%D0\%B0\%D1\%87_\%D0\%B8_\%D0\%BF\%D0\%BB\%D0\%B0\%D0\%BD\%D0\%B8\%D1\%80\%D0\%BE\%D0\%B2\%D1\%89\%D0\%B8\%D0\%BA.pdf\#nameddest=swr-ses-021-\%D0\%B7\%D0\%B0\%D0\%BF\%D0\%B8\%D1\%81\%D1\%8C-\%D0\%BF\%D1\%80\%D0\%B8\%D0\%BD\%D1\%8F\%D1\%82\%D0\%BE\%D0\%B9-\%D0\%BE\%D0\%BF\%D0\%B5\%D1\%80\%D0\%B0\%D1\%86\%D0\%B8\%D0\%B8}{сохранённую запись операции}.

Источник: SR-SYS-005. Проверка: сбой сохранения не оставляет подтверждённый пользователю вердикт с противоречащим ему состоянием сессии.

\subsection{SWR-SYS-006. Общий серверный срок}\label{swr-sys-006-ux43eux431ux449ux438ux439-ux441ux435ux440ux432ux435ux440ux43dux44bux439-ux441ux440ux43eux43a}

\hypertarget{swr-sys-006-общий-серверный-срок}{}

Программные части не должны получать разные интервалы решения из собственных таймеров. Общий срок выводится из сохранённого времени создания сессии; своевременность отдельного запроса --- из его времени поступления. Эти значения нужны и обработчику ответа, и фоновому контролю, и восстановлению после перезапуска.

\textbf{Требование.} Обработка ответов и контроль просрочки должны использовать одно сохранённое время создания сессии, один срок и зафиксированное сервером время получения ответа.

После перезапуска начало работы приложения не заменяет начало набора. После ожидания в очереди старт проверяющего не заменяет момент приёма. Отдельный предел обработки тоже продолжает отсчитываться от прежнего \texttt{t\_in}. Обозначения даны в \href{../03_\%D0\%9A\%D0\%B0\%D1\%80\%D1\%82\%D0\%BE\%D1\%87\%D0\%BA\%D0\%B8_\%D0\%B4\%D0\%BE\%D0\%BC\%D0\%B5\%D0\%BD\%D0\%BE\%D0\%B2/\%D0\%A1\%D0\%B5\%D1\%81\%D1\%81\%D0\%B8\%D0\%B8_\%D0\%B7\%D0\%B0\%D0\%B4\%D0\%B0\%D1\%87.pdf\#nameddest=\%D1\%81\%D1\%80\%D0\%BE\%D0\%BA-\%D0\%B8-\%D0\%BF\%D0\%BE\%D0\%BF\%D1\%8B\%D1\%82\%D0\%BA\%D0\%B0}{карточке домена}, а программные обязанности разделены между \href{DOC-06.2_\%D0\%A1\%D0\%B5\%D1\%81\%D1\%81\%D0\%B8\%D0\%B8_\%D0\%B7\%D0\%B0\%D0\%B4\%D0\%B0\%D1\%87_\%D0\%B8_\%D0\%BF\%D0\%BB\%D0\%B0\%D0\%BD\%D0\%B8\%D1\%80\%D0\%BE\%D0\%B2\%D1\%89\%D0\%B8\%D0\%BA.pdf\#nameddest=swr-ses-006-\%D1\%81\%D0\%BE\%D1\%85\%D1\%80\%D0\%B0\%D0\%BD\%D0\%B5\%D0\%BD\%D0\%B8\%D0\%B5-\%D0\%B2\%D1\%80\%D0\%B5\%D0\%BC\%D0\%B5\%D0\%BD\%D0\%B8-\%D1\%81\%D0\%BE\%D0\%B7\%D0\%B4\%D0\%B0\%D0\%BD\%D0\%B8\%D1\%8F}{сохранением времени}, \href{DOC-06.2_\%D0\%A1\%D0\%B5\%D1\%81\%D1\%81\%D0\%B8\%D0\%B8_\%D0\%B7\%D0\%B0\%D0\%B4\%D0\%B0\%D1\%87_\%D0\%B8_\%D0\%BF\%D0\%BB\%D0\%B0\%D0\%BD\%D0\%B8\%D1\%80\%D0\%BE\%D0\%B2\%D1\%89\%D0\%B8\%D0\%BA.pdf\#nameddest=swr-ses-007-\%D1\%80\%D0\%B5\%D0\%B3\%D0\%B8\%D1\%81\%D1\%82\%D1\%80\%D0\%B0\%D1\%86\%D0\%B8\%D1\%8F-\%D0\%B2\%D1\%80\%D0\%B5\%D0\%BC\%D0\%B5\%D0\%BD\%D0\%B8-\%D0\%BF\%D0\%BE\%D1\%81\%D1\%82\%D1\%83\%D0\%BF\%D0\%BB\%D0\%B5\%D0\%BD\%D0\%B8\%D1\%8F}{регистрацией поступления} и \href{DOC-06.2_\%D0\%A1\%D0\%B5\%D1\%81\%D1\%81\%D0\%B8\%D0\%B8_\%D0\%B7\%D0\%B0\%D0\%B4\%D0\%B0\%D1\%87_\%D0\%B8_\%D0\%BF\%D0\%BB\%D0\%B0\%D0\%BD\%D0\%B8\%D1\%80\%D0\%BE\%D0\%B2\%D1\%89\%D0\%B8\%D0\%BA.pdf\#nameddest=swr-ses-020-\%D0\%BF\%D1\%80\%D0\%B5\%D0\%B4\%D0\%B5\%D0\%BB-\%D0\%BE\%D0\%B1\%D1\%80\%D0\%B0\%D0\%B1\%D0\%BE\%D1\%82\%D0\%BA\%D0\%B8}{пределом обработки}.

Источники: SR-SYS-006, SR-SYS-007. Проверка: перезапуск не продлевает срок, а задержка планировщика не меняет допустимость ответа.

\subsection{SWR-SYS-007. Категория результата операции}\label{swr-sys-007-ux43aux430ux442ux435ux433ux43eux440ux438ux44f-ux440ux435ux437ux443ux43bux44cux442ux430ux442ux430-ux43eux43fux435ux440ux430ux446ux438ux438}

\hypertarget{swr-sys-007-категория-результата-операции}{}

Внешний интерфейс получает результаты от нескольких программных частей. Он должен сохранить их категорию, чтобы клиент отличал предметное решение от отказа выполнить операцию. Самого текста сообщения для этого недостаточно: одинаковые слова не определяют ни состояние задачи, ни возможность повторного ответа.

\textbf{Требование.} Компонент обработки запроса должен сохранять различие между предметным вердиктом и отказами при передаче результата внешнему интерфейсу.

Сохранённый неверный вердикт передаётся предметным результатом с \texttt{isOk\ =\ false}. Ошибка формата, нарушение доступа, просрочка и технический сбой передаются машинными кодами из \href{DOC-06.2_\%D0\%A1\%D0\%B5\%D1\%81\%D1\%81\%D0\%B8\%D0\%B8_\%D0\%B7\%D0\%B0\%D0\%B4\%D0\%B0\%D1\%87_\%D0\%B8_\%D0\%BF\%D0\%BB\%D0\%B0\%D0\%BD\%D0\%B8\%D1\%80\%D0\%BE\%D0\%B2\%D1\%89\%D0\%B8\%D0\%BA.pdf\#nameddest=\%D0\%BA\%D0\%BE\%D0\%B4\%D1\%8B-\%D0\%B8-\%D0\%BF\%D1\%80\%D0\%B5\%D0\%B4\%D1\%81\%D1\%82\%D0\%B0\%D0\%B2\%D0\%BB\%D0\%B5\%D0\%BD\%D0\%B8\%D0\%B5-\%D0\%BE\%D1\%88\%D0\%B8\%D0\%B1\%D0\%BE\%D0\%BA}{контракта ошибок}. При преобразовании внутренних результатов нельзя превращать исключение проверяющего в отрицательный предметный ответ. Эти различия следуют из \href{../04_\%D0\%A1\%D0\%B8\%D1\%81\%D1\%82\%D0\%B5\%D0\%BC\%D0\%BD\%D1\%8B\%D0\%B5_\%D1\%82\%D1\%80\%D0\%B5\%D0\%B1\%D0\%BE\%D0\%B2\%D0\%B0\%D0\%BD\%D0\%B8\%D1\%8F/DOC-05_\%D0\%A1\%D0\%B8\%D1\%81\%D1\%82\%D0\%B5\%D0\%BC\%D0\%BD\%D1\%8B\%D0\%B5_\%D1\%82\%D1\%80\%D0\%B5\%D0\%B1\%D0\%BE\%D0\%B2\%D0\%B0\%D0\%BD\%D0\%B8\%D1\%8F.pdf\#nameddest=sr-sys-009-\%D1\%80\%D0\%B0\%D0\%B7\%D0\%BB\%D0\%B8\%D1\%87\%D0\%B5\%D0\%BD\%D0\%B8\%D0\%B5-\%D0\%BE\%D1\%82\%D0\%BA\%D0\%B0\%D0\%B7\%D0\%BE\%D0\%B2}{системного требования категорий отказа}.

Источники: SR-SYS-008, SR-SYS-009. Проверка: неверное решение, нарушение доступа, ошибка формата, просрочка и технический сбой распознаются как разные категории; соответствие HTTP-кодам задано контрактом сессий и сохраняется в общем справочнике DOC-08.

\subsection{SWR-SYS-008. Сохранность завершённых данных}\label{swr-sys-008-ux441ux43eux445ux440ux430ux43dux43dux43eux441ux442ux44c-ux437ux430ux432ux435ux440ux448ux451ux43dux43dux44bux445-ux434ux430ux43dux43dux44bux445}

\hypertarget{swr-sys-008-сохранность-завершённых-данных}{}

Состояние системы должно переживать завершение процесса приложения. В хранилище остаются не только конечные значения, но и принадлежность, условия и временные отметки, необходимые для их чтения. Иначе после восстановления нельзя было бы определить, кому принадлежит результат и по какому заданию он получен.

\textbf{Требование.} Хранение должно обеспечивать восстановление идентификаторов, принадлежности, условий, времени создания и завершённых результатов сессий после перезапуска приложения.

Готовый вердикт используется как сохранённый факт операции, без новой математической попытки. Принятая операция без вердикта восстанавливается по собственной записи, сохраняя предел времени. Это различие связывает \href{../04_\%D0\%A1\%D0\%B8\%D1\%81\%D1\%82\%D0\%B5\%D0\%BC\%D0\%BD\%D1\%8B\%D0\%B5_\%D1\%82\%D1\%80\%D0\%B5\%D0\%B1\%D0\%BE\%D0\%B2\%D0\%B0\%D0\%BD\%D0\%B8\%D1\%8F/DOC-05_\%D0\%A1\%D0\%B8\%D1\%81\%D1\%82\%D0\%B5\%D0\%BC\%D0\%BD\%D1\%8B\%D0\%B5_\%D1\%82\%D1\%80\%D0\%B5\%D0\%B1\%D0\%BE\%D0\%B2\%D0\%B0\%D0\%BD\%D0\%B8\%D1\%8F.pdf\#nameddest=sr-sys-007-\%D1\%81\%D0\%BE\%D1\%85\%D1\%80\%D0\%B0\%D0\%BD\%D0\%BD\%D0\%BE\%D1\%81\%D1\%82\%D1\%8C-\%D0\%B4\%D0\%B0\%D0\%BD\%D0\%BD\%D1\%8B\%D1\%85-\%D1\%81\%D0\%B5\%D1\%81\%D1\%81\%D0\%B8\%D0\%B8}{системную сохранность} с \href{DOC-06.2_\%D0\%A1\%D0\%B5\%D1\%81\%D1\%81\%D0\%B8\%D0\%B8_\%D0\%B7\%D0\%B0\%D0\%B4\%D0\%B0\%D1\%87_\%D0\%B8_\%D0\%BF\%D0\%BB\%D0\%B0\%D0\%BD\%D0\%B8\%D1\%80\%D0\%BE\%D0\%B2\%D1\%89\%D0\%B8\%D0\%BA.pdf\#nameddest=swr-ses-021-\%D0\%B7\%D0\%B0\%D0\%BF\%D0\%B8\%D1\%81\%D1\%8C-\%D0\%BF\%D1\%80\%D0\%B8\%D0\%BD\%D1\%8F\%D1\%82\%D0\%BE\%D0\%B9-\%D0\%BE\%D0\%BF\%D0\%B5\%D1\%80\%D0\%B0\%D1\%86\%D0\%B8\%D0\%B8}{восстановлением операции} и \href{DOC-06.2_\%D0\%A1\%D0\%B5\%D1\%81\%D1\%81\%D0\%B8\%D0\%B8_\%D0\%B7\%D0\%B0\%D0\%B4\%D0\%B0\%D1\%87_\%D0\%B8_\%D0\%BF\%D0\%BB\%D0\%B0\%D0\%BD\%D0\%B8\%D1\%80\%D0\%BE\%D0\%B2\%D1\%89\%D0\%B8\%D0\%BA.pdf\#nameddest=swr-ses-023-\%D1\%87\%D1\%82\%D0\%B5\%D0\%BD\%D0\%B8\%D0\%B5-\%D1\%81\%D0\%B5\%D1\%81\%D1\%81\%D0\%B8\%D0\%B8-\%D0\%B8-\%D1\%80\%D0\%B5\%D0\%B7\%D1\%83\%D0\%BB\%D1\%8C\%D1\%82\%D0\%B0\%D1\%82\%D0\%B0}{чтением сессии}. Подтверждённый результат должен быть доступен клиенту даже после потери исходного соединения.

Источник: SR-SYS-007. Проверка: после восстановления выдаётся прежний итог и применяется прежний срок. Для принятого незавершённого ответа используются сохранённая операция и прежний срок обработки по DOC-06.2.

\section{Состав подсистемных документов}\label{4-ux441ux43eux441ux442ux430ux432-ux43fux43eux434ux441ux438ux441ux442ux435ux43cux43dux44bux445-ux434ux43eux43aux443ux43cux435ux43dux442ux43eux432}

\hypertarget{4-состав-подсистемных-документов}{}

{\def\LTcaptype{none} % do not increment counter
\begin{longtable}[]{@{}
  >{\raggedright\arraybackslash}p{(\linewidth - 4\tabcolsep) * \real{0.3276}}
  >{\raggedright\arraybackslash}p{(\linewidth - 4\tabcolsep) * \real{0.3276}}
  >{\raggedright\arraybackslash}p{(\linewidth - 4\tabcolsep) * \real{0.3447}}@{}}
\toprule\noalign{}
\begin{minipage}[b]{\linewidth}\raggedright
Документ
\end{minipage} & \begin{minipage}[b]{\linewidth}\raggedright
Владелец
\end{minipage} & \begin{minipage}[b]{\linewidth}\raggedright
Предметная область
\end{minipage} \\
\midrule\noalign{}
\endhead
\bottomrule\noalign{}
\endlastfoot
DOC-06.1 & Кокорев & Аутентификация и пользователи \\
\href{DOC-06.2_\%D0\%A1\%D0\%B5\%D1\%81\%D1\%81\%D0\%B8\%D0\%B8_\%D0\%B7\%D0\%B0\%D0\%B4\%D0\%B0\%D1\%87_\%D0\%B8_\%D0\%BF\%D0\%BB\%D0\%B0\%D0\%BD\%D0\%B8\%D1\%80\%D0\%BE\%D0\%B2\%D1\%89\%D0\%B8\%D0\%BA.pdf\#nameddest=doc-062-\%D1\%81\%D0\%B5\%D1\%81\%D1\%81\%D0\%B8\%D0\%B8-\%D0\%B7\%D0\%B0\%D0\%B4\%D0\%B0\%D1\%87-\%D0\%B8-\%D0\%BF\%D0\%BB\%D0\%B0\%D0\%BD\%D0\%B8\%D1\%80\%D0\%BE\%D0\%B2\%D1\%89\%D0\%B8\%D0\%BA}{DOC-06.2} & Анохин & Сессии задач и планировщик \\
DOC-06.3 & Георгиев & Ранцевые задачи \\
DOC-06.4 & Минасов & Учебный бэкдор-сценарий \\
DOC-06.5 & Портнов & Эль-Гамаль \\
DOC-06.6 & Самбьюи & Конфигурации и модерация \\
\end{longtable}
}

\subsection{Обязательства на границах подсистем}\label{ux43eux431ux44fux437ux430ux442ux435ux43bux44cux441ux442ux432ux430-ux43dux430-ux433ux440ux430ux43dux438ux446ux430ux445-ux43fux43eux434ux441ux438ux441ux442ux435ux43c}

\hypertarget{обязательства-на-границах-подсистем}{}

Безопасность задаёт подтверждённого пользователя и условия доступа; сессии используют этот результат. Конфигурации предоставляют доступный состав, по которому создаётся набор. Предметные обработчики различают вердикт и невозможность проверки; общий интерфейс сохраняет это различие при сообщении результата.

Задача, итог сессии и запись обработанного ответа сохраняются в одной транзакции. Незавершённая запись служит основанием восстановления. Время получения ответа, общий срок и предел обработки сохраняются при повторном запуске; вычислительная задержка не превращает своевременный ответ в поздний.

\section{Связи и проверка редакции}\label{5-ux441ux432ux44fux437ux438-ux438-ux43fux440ux43eux432ux435ux440ux43aux430-ux440ux435ux434ux430ux43aux446ux438ux438}

\hypertarget{5-связи-и-проверка-редакции}{}

Требования SR-SYS-001--SR-SYS-009 раскрыты восемью SWR-SYS выше. Для каждого указаны источник и ожидаемая проверка. Детализация SR-SES находится в DOC-06.2. Описательные кейсы DOC-10 связываются с этими требованиями в матрице DOC-09.

Проверка редакции охватывает единый контракт ошибок, точность целых чисел, доступ по владельцу, согласованность сохранения и восстановление принятого ответа. \href{../04_\%D0\%A1\%D0\%B8\%D1\%81\%D1\%82\%D0\%B5\%D0\%BC\%D0\%BD\%D1\%8B\%D0\%B5_\%D1\%82\%D1\%80\%D0\%B5\%D0\%B1\%D0\%BE\%D0\%B2\%D0\%B0\%D0\%BD\%D0\%B8\%D1\%8F/DOC-05_\%D0\%A1\%D0\%B8\%D1\%81\%D1\%82\%D0\%B5\%D0\%BC\%D0\%BD\%D1\%8B\%D0\%B5_\%D1\%82\%D1\%80\%D0\%B5\%D0\%B1\%D0\%BE\%D0\%B2\%D0\%B0\%D0\%BD\%D0\%B8\%D1\%8F.pdf\#nameddest=5-\%D1\%80\%D0\%B0\%D1\%81\%D0\%BF\%D1\%80\%D0\%B5\%D0\%B4\%D0\%B5\%D0\%BB\%D0\%B5\%D0\%BD\%D0\%B8\%D0\%B5-\%D0\%BE\%D1\%82\%D0\%B2\%D0\%B5\%D1\%82\%D1\%81\%D1\%82\%D0\%B2\%D0\%B5\%D0\%BD\%D0\%BD\%D0\%BE\%D1\%81\%D1\%82\%D0\%B8}{DOC-05} задаёт распределение ответственности.

История: 05.11.2026 --- подготовлена версия 0.5: ответственность частей, общие программные обязательства, связи с системными требованиями и состав подсистем. 07.11.2026: расширено описание взаимодействий и определено восстановление; подготовлены LaTeX и PDF.

Редакция 07.11.2026: развёрнуты пояснения каждого требования, добавлены примеры и точные ссылки на связанные разделы. Формулировки обязательств и идентификаторы сохранены.


\end{document}
